\documentclass[10pt,a4paper]{amsart}
\usepackage{amsmath,amssymb,amsthm,graphicx}
\usepackage{verbatim}
\newenvironment{reference}{\comment}{\endcomment}
%\newenvironment{reference}{}{}
\newenvironment{slogan}{\comment}{\endcomment}
\newenvironment{history}{\comment}{\endcomment}

% For commutative diagrams we use Xy-pic
\usepackage[all]{xy}

% We use 2cell for 2-commutative diagrams.
\xyoption{2cell}
% Theorem environments.
%
\theoremstyle{plain}
\newtheorem{theorem}[subsection]{Theorem}
\newtheorem{proposition}[subsection]{Proposition}
\newtheorem{lemma}[subsection]{Lemma}

\theoremstyle{definition}
\newtheorem{definition}[subsection]{Definition}
\newtheorem{example}[subsection]{Example}
\newtheorem{exercise}[subsection]{Exercise}
\newtheorem{situation}[subsection]{Situation}

\theoremstyle{remark}
\newtheorem{remark}[subsection]{Remark}
\newtheorem{remarks}[subsection]{Remarks}

\numberwithin{equation}{subsection}

% Macros
%
\def\lim{\mathop{\mathrm{lim}}\nolimits}
\def\colim{\mathop{\mathrm{colim}}\nolimits}
\def\Spec{\mathop{\mathrm{Spec}}}
\def\Hom{\mathop{\mathrm{Hom}}\nolimits}
\def\Ext{\mathop{\mathrm{Ext}}\nolimits}
\def\SheafHom{\mathop{\mathcal{H}\!\mathit{om}}\nolimits}
\def\SheafExt{\mathop{\mathcal{E}\!\mathit{xt}}\nolimits}
\def\Sch{\mathit{Sch}}
\def\Mor{\mathop{\mathrm{Mor}}\nolimits}
\def\Ob{\mathop{\mathrm{Ob}}\nolimits}
\def\Sh{\mathop{\mathit{Sh}}\nolimits}
\def\NL{\mathop{N\!L}\nolimits}
\def\CH{\mathop{\mathrm{CH}}\nolimits}
\def\proetale{{pro\text{-}\acute{e}tale}}
\def\etale{{\acute{e}tale}}
\def\QCoh{\mathit{QCoh}}
\def\Ker{\mathop{\mathrm{Ker}}}
\def\Im{\mathop{\mathrm{Im}}}
\def\Coker{\mathop{\mathrm{Coker}}}
\def\Coim{\mathop{\mathrm{Coim}}}

% Boxtimes
%
\DeclareMathSymbol{\boxtimes}{\mathbin}{AMSa}{"02}

%
% Macros for moduli stacks/spaces
%
\def\QCohstack{\mathcal{QC}\!\mathit{oh}}
\def\Cohstack{\mathcal{C}\!\mathit{oh}}
\def\Spacesstack{\mathcal{S}\!\mathit{paces}}
\def\Quotfunctor{\mathrm{Quot}}
\def\Hilbfunctor{\mathrm{Hilb}}
\def\Curvesstack{\mathcal{C}\!\mathit{urves}}
\def\Polarizedstack{\mathcal{P}\!\mathit{olarized}}
\def\Complexesstack{\mathcal{C}\!\mathit{omplexes}}
% \Pic is the operator that assigns to X its picard group, usage \Pic(X)
% \Picardstack_{X/B} denotes the Picard stack of X over B
% \Picardfunctor_{X/B} denotes the Picard functor of X over B
\def\Pic{\mathop{\mathrm{Pic}}\nolimits}
\def\Picardstack{\mathcal{P}\!\mathit{ic}}
\def\Picardfunctor{\mathrm{Pic}}
\def\Deformationcategory{\mathcal{D}\!\mathit{ef}}
\usepackage{cleveref}
\crefname{originallemma}{original lemma}{original lemmas}
\Crefname{originallemma}{Original lemma}{Original lemmas}
\makeatletter
\newlabel{cohomology-lemma-h1-extensions}{{0B3A}{source}{}{sourcepointer-1}{}}
\newlabel{cohomology-lemma-h1-extensions@cref}{{[originallemma][0][]0B3A}{[1][0][]source}{}{}{}}
\makeatother
\crefname{originallemma}{original lemma}{original lemmas}
\Crefname{originallemma}{Original lemma}{Original lemmas}
\makeatletter
\newlabel{divisors-lemma-conormal-sheaf-section-projective-bundle}{{0B3V}{source}{}{sourcepointer-2}{}}
\newlabel{divisors-lemma-conormal-sheaf-section-projective-bundle@cref}{{[originallemma][0][]0B3V}{[1][0][]source}{}{}{}}
\makeatother
\crefname{originallemma}{original lemma}{original lemmas}
\Crefname{originallemma}{Original lemma}{Original lemmas}
\makeatletter
\newlabel{divisors-lemma-section-smooth-regular-immersion}{{067R}{source}{}{sourcepointer-3}{}}
\newlabel{divisors-lemma-section-smooth-regular-immersion@cref}{{[originallemma][0][]067R}{[1][0][]source}{}{}{}}
\makeatother
\crefname{originalreference}{original reference}{original references}
\Crefname{originalreference}{Original reference}{Original references}
\makeatletter
\newlabel{lemma-nontrivial-normal-bundle}{{0AXA}{source}{}{sourcepointer-4}{}}
\newlabel{lemma-nontrivial-normal-bundle@cref}{{[originalreference][0][]0AXA}{[1][0][]source}{}{}{}}
\makeatother
\crefname{originaldefinition}{original definition}{original definitions}
\Crefname{originaldefinition}{Original definition}{Original definitions}
\makeatletter
\newlabel{morphisms-definition-scheme-theoretic-image}{{01R7}{source}{}{sourcepointer-5}{}}
\newlabel{morphisms-definition-scheme-theoretic-image@cref}{{[originaldefinition][0][]01R7}{[1][0][]source}{}{}{}}
\makeatother
\crefname{originallemma}{original lemma}{original lemmas}
\Crefname{originallemma}{Original lemma}{Original lemmas}
\makeatletter
\newlabel{morphisms-lemma-conormal-functorial-flat}{{0473}{source}{}{sourcepointer-6}{}}
\newlabel{morphisms-lemma-conormal-functorial-flat@cref}{{[originallemma][0][]0473}{[1][0][]source}{}{}{}}
\makeatother
\crefname{originallemma}{original lemma}{original lemmas}
\Crefname{originallemma}{Original lemma}{Original lemmas}
\makeatletter
\newlabel{morphisms-lemma-finite-birational-over-normal}{{0AB1}{source}{}{sourcepointer-7}{}}
\newlabel{morphisms-lemma-finite-birational-over-normal@cref}{{[originallemma][0][]0AB1}{[1][0][]source}{}{}{}}
\makeatother
\crefname{originallemma}{original lemma}{original lemmas}
\Crefname{originallemma}{Original lemma}{Original lemmas}
\makeatletter
\newlabel{morphisms-lemma-finite-proper}{{01WN}{source}{}{sourcepointer-8}{}}
\newlabel{morphisms-lemma-finite-proper@cref}{{[originallemma][0][]01WN}{[1][0][]source}{}{}{}}
\makeatother
\crefname{originallemma}{original lemma}{original lemmas}
\Crefname{originallemma}{Original lemma}{Original lemmas}
\makeatletter
\newlabel{morphisms-lemma-nagata-normalization}{{035S}{source}{}{sourcepointer-9}{}}
\newlabel{morphisms-lemma-nagata-normalization@cref}{{[originallemma][0][]035S}{[1][0][]source}{}{}{}}
\makeatother
\crefname{originallemma}{original lemma}{original lemmas}
\Crefname{originallemma}{Original lemma}{Original lemmas}
\makeatletter
\newlabel{morphisms-lemma-ubiquity-nagata}{{035B}{source}{}{sourcepointer-10}{}}
\newlabel{morphisms-lemma-ubiquity-nagata@cref}{{[originallemma][0][]035B}{[1][0][]source}{}{}{}}
\makeatother
\crefname{originalreference}{original reference}{original references}
\Crefname{originalreference}{Original reference}{Original references}
\makeatletter
\newlabel{situation-vanishing}{{07UY}{source}{}{sourcepointer-11}{}}
\newlabel{situation-vanishing@cref}{{[originalreference][0][]07UY}{[1][0][]source}{}{}{}}
\makeatother
\title{Injectivity of the restriction of the first cohomology of the structure sheaf to the inverse image of the open}
\author{The Stacks Project Authors}
\date{stacks-project:resolve.tex:local-snapshot-20261012; GNU FDL 1.2 or later}
\begin{document}
\maketitle
\noindent\textit{Editorial note.} All delivered mathematical text is the author's own, copied from the pinned Stacks Project chapter file resolve.tex; no statement or proof step was rewritten, reordered or removed.\par\medskip
\noindent\textit{Editorial note.} The counted proof is the author's complete argument: the translation of a nontrivial class into an extension with its projective bundle and the second section produced by triviality of the restriction, the argument that the two sections must meet through affineness and finiteness of the scheme theoretic image, and the normalization step with the computation of the conormal sheaf of the pullback divisor leading to the contradiction.\par\medskip
\noindent\textit{Editorial note.} The label lemma-H1-injective names the counted statement, and the proof environment immediately following it in the source is delivered verbatim as its proof.\par\medskip
\noindent\textit{Editorial note.} The counted statement is carried by the author's own lemma environment in the chapter source.\par\medskip
\noindent\textit{Editorial note.} Pre-open gate: python operations/agent-authoring/tools/precard\_gate.py --source resolve.tex --label lemma-H1-injective --start 1691 --end 1754 --local-id 28077 returned RESULT=OK (L28077 free, occupancy FREE, public interval clear). Pre-assembly tag self-check: canonical tag resolved and not present in reference/public-items.jsonl.\par\medskip
\noindent\textit{Editorial note.} The statement and proof are copied from the pinned chapter file; the cited inputs are kept as external pointers with their canonical Stacks tags.\par\medskip
\noindent\textit{Allowed standard prerequisites (editorial).}\begin{itemize}
\item Classification of classes by extensions and the projective bundle of the extension. Cohomology, the lemma identifying first cohomology classes with extensions of the structure sheaf, and Divisors, the lemma describing the conormal sheaf of the section of a projective bundle defined by a surjection. Supplies the extension corresponding to a nontrivial class, its projective bundle and the section whose conormal sheaf is the structure sheaf.
\item Affineness of the complement of the section and finiteness of an affine proper morphism. Standard properties of projective bundles: the complement of the image of a section is affine; Morphisms, the lemma that an affine and proper morphism is finite, and the lemma that a finite birational morphism onto a normal scheme is an isomorphism. Supply the contradiction argument showing that the two sections must meet, and the normality step for the scheme theoretic image.
\item Finiteness of normalization for Nagata schemes and Cartier divisors from sections. Morphisms, the lemnas that normalization of a Nagata scheme is finite and that the image of a section of a smooth morphism is an effective Cartier divisor, and Divisors, the consequence that the pullback to the normalization is again a Cartier divisor. Supply the normalized scheme, the Cartier divisor obtained by pulling back the section, and the computation of the conormal sheaf on the normalization.
\item Impossibility of a nontrivial invertible conormal sheaf on such a divisor. Same chapter, the lemma asserting the nontriviality obstruction for the conormal sheaf of the Cartier divisor obtained above. Supplies the final contradiction, which proves the injectivity.
\end{itemize}\medskip
\ifdefined\hypertarget\hypertarget{sourcepointer-1}{}\fi
\noindent\textit{Source reference.} Original lemma 0B3A: sources/stacks-chow/tags.txt, source lines 13958--13958. Prerequisite pointer invoked at proof line(s) 1705; it names a standard fact of a sibling chapter that the counted argument uses, and every step of the counted proof built on it is retained verbatim in the delivered body.\par\medskip
\ifdefined\hypertarget\hypertarget{sourcepointer-2}{}\fi
\noindent\textit{Source reference.} Original lemma 0B3V: sources/stacks-chow/tags.txt, source lines 13978--13978. Prerequisite pointer invoked at proof line(s) 1711; it names a standard fact of a sibling chapter that the counted argument uses, and every step of the counted proof built on it is retained verbatim in the delivered body.\par\medskip
\ifdefined\hypertarget\hypertarget{sourcepointer-3}{}\fi
\noindent\textit{Source reference.} Original lemma 067R: sources/stacks-chow/tags.txt, source lines 7980--7980. Prerequisite pointer invoked at proof line(s) 1745; it names a standard fact of a sibling chapter that the counted argument uses, and every step of the counted proof built on it is retained verbatim in the delivered body.\par\medskip
\ifdefined\hypertarget\hypertarget{sourcepointer-4}{}\fi
\noindent\textit{Source reference.} Original reference 0AXA: sources/stacks-resolve/pinned-original-source.tex, source lines 1644--1644. Prerequisite pointer invoked at proof line(s) 1752; it names a standard fact of the same chapter that the counted argument uses, and every step of the counted proof built on it is retained verbatim in the delivered body.\par\medskip
\ifdefined\hypertarget\hypertarget{sourcepointer-5}{}\fi
\noindent\textit{Source reference.} Original definition 01R7: sources/stacks-chow/tags.txt, source lines 2292--2292. Prerequisite pointer invoked at proof line(s) 1719; it names a standard fact of a sibling chapter that the counted argument uses, and every step of the counted proof built on it is retained verbatim in the delivered body.\par\medskip
\ifdefined\hypertarget\hypertarget{sourcepointer-6}{}\fi
\noindent\textit{Source reference.} Original lemma 0473: sources/stacks-chow/tags.txt, source lines 5494--5494. Prerequisite pointer invoked at proof line(s) 1750; it names a standard fact of a sibling chapter that the counted argument uses, and every step of the counted proof built on it is retained verbatim in the delivered body.\par\medskip
\ifdefined\hypertarget\hypertarget{sourcepointer-7}{}\fi
\noindent\textit{Source reference.} Original lemma 0AB1: sources/stacks-chow/tags.txt, source lines 13004--13004. Prerequisite pointer invoked at proof line(s) 1732; it names a standard fact of a sibling chapter that the counted argument uses, and every step of the counted proof built on it is retained verbatim in the delivered body.\par\medskip
\ifdefined\hypertarget\hypertarget{sourcepointer-8}{}\fi
\noindent\textit{Source reference.} Original lemma 01WN: sources/stacks-chow/tags.txt, source lines 2488--2488. Prerequisite pointer invoked at proof line(s) 1730; it names a standard fact of a sibling chapter that the counted argument uses, and every step of the counted proof built on it is retained verbatim in the delivered body.\par\medskip
\ifdefined\hypertarget\hypertarget{sourcepointer-9}{}\fi
\noindent\textit{Source reference.} Original lemma 035S: sources/stacks-chow/tags.txt, source lines 4139--4139. Prerequisite pointer invoked at proof line(s) 1738; it names a standard fact of a sibling chapter that the counted argument uses, and every step of the counted proof built on it is retained verbatim in the delivered body.\par\medskip
\ifdefined\hypertarget\hypertarget{sourcepointer-10}{}\fi
\noindent\textit{Source reference.} Original lemma 035B: sources/stacks-chow/tags.txt, source lines 4122--4122. Prerequisite pointer invoked at proof line(s) 1739; it names a standard fact of a sibling chapter that the counted argument uses, and every step of the counted proof built on it is retained verbatim in the delivered body.\par\medskip
\ifdefined\hypertarget\hypertarget{sourcepointer-11}{}\fi
\noindent\textit{Source reference.} Original reference 07UY: sources/stacks-resolve/pinned-original-source.tex, source lines 1584--1584. Prerequisite pointer invoked at proof line(s) 1693; it names a standard fact of the same chapter that the counted argument uses, and every step of the counted proof built on it is retained verbatim in the delivered body.\par\medskip
\begin{lemma}
\label{lemma-H1-injective}
In Situation \ref{situation-vanishing} assume $X$ is normal
and $A$ Nagata. The map
$$
H^1(X, \mathcal{O}_X) \longrightarrow H^1(f^{-1}(U), \mathcal{O}_X)
$$
is injective.
\end{lemma}

\begin{proof}
Let $0 \to \mathcal{O}_X \to \mathcal{E} \to \mathcal{O}_X \to 0$ be the
extension corresponding to a nontrivial element $\xi$ of
$H^1(X, \mathcal{O}_X)$
(Cohomology, Lemma \ref{cohomology-lemma-h1-extensions}).
Let $\pi : P = \mathbf{P}(\mathcal{E}) \to X$
be the projective bundle associated to $\mathcal{E}$.
The surjection $\mathcal{E} \to \mathcal{O}_X$
defines a section $\sigma : X \to P$ whose conormal sheaf is
isomorphic to $\mathcal{O}_X$ (Divisors, Lemma
\ref{divisors-lemma-conormal-sheaf-section-projective-bundle}).
If the restriction of $\xi$ to $f^{-1}(U)$ is trivial, then we get
a map $\mathcal{E}|_{f^{-1}(U)} \to \mathcal{O}_{f^{-1}(U)}$ splitting
the injection $\mathcal{O}_X \to \mathcal{E}$. This defines a second
section $\sigma' : f^{-1}(U) \to P$ disjoint from $\sigma$. Since $\xi$
is nontrivial we conclude that $\sigma'$ cannot extend to all of $X$
and be disjoint from $\sigma$. Let $X' \subset P$ be the
scheme theoretic image of $\sigma'$ (Morphisms,
Definition \ref{morphisms-definition-scheme-theoretic-image}).
Picture
$$
\xymatrix{
& X' \ar[rd]_g \ar[r] & P \ar[d]_\pi \\
f^{-1}(U) \ar[ru]_{\sigma'} \ar[rr] & & X \ar@/_/[u]_\sigma
}
$$
The morphism $P \setminus \sigma(X) \to X$ is affine.
If $X' \cap \sigma(X) = \emptyset$, then $X' \to X$ is both affine
and proper, hence finite
(Morphisms, Lemma \ref{morphisms-lemma-finite-proper}),
hence an isomorphism (as $X$ is normal, see
Morphisms, Lemma \ref{morphisms-lemma-finite-birational-over-normal}).
This is impossible as mentioned above.

\medskip\noindent
Let $X^\nu$ be the normalization of $X'$.
Since $A$ is Nagata, we see that $X^\nu \to X'$ is finite
(Morphisms, Lemmas \ref{morphisms-lemma-nagata-normalization} and
\ref{morphisms-lemma-ubiquity-nagata}). Let $Z \subset X^\nu$ be the
pullback of the effective Cartier divisor $\sigma(X) \subset P$.
By the above we see that $Z$ is not empty and is contained
in the closed fibre of $X^\nu \to S$.
Since $P \to X$ is smooth, we see that $\sigma(X)$ is an effective
Cartier divisor
(Divisors, Lemma \ref{divisors-lemma-section-smooth-regular-immersion}).
Hence $Z \subset X^\nu$ is an effective Cartier divisor too.
Since the conormal sheaf of $\sigma(X)$ in $P$ is $\mathcal{O}_X$, the
conormal sheaf of $Z$ in $X^\nu$ (which is a priori invertible)
is $\mathcal{O}_Z$ by
Morphisms, Lemma \ref{morphisms-lemma-conormal-functorial-flat}.
This is impossible by
Lemma \ref{lemma-nontrivial-normal-bundle}
and the proof is complete.
\end{proof}

\end{document}
